\documentclass{article}
\usepackage[margin=1in]{geometry}
\usepackage{enumitem}
\usepackage{listings}

\title{Checkpoint 5 Test Notes}
\author{}
\date{}

\begin{document}
\maketitle

\section*{Scope}
Checkpoint 5 requires complete \texttt{TtyRead} and \texttt{TtyWrite}
syscalls, receive/transmit terminal trap handling, and the remaining trap
handlers.

\section*{Checks}
\begin{itemize}[leftmargin=*]
  \item \texttt{make clean \&\& make}: clean build succeeds with
        \texttt{YALNIX\_FRAMEWORK} set.
  \item \texttt{cp4\_fork}, \texttt{cp4\_exec}, and \texttt{cp4\_stack}: all
        exit with status 0 under \texttt{./yalnix -W}.
  \item \texttt{cp5\_tty}: writes a short message and then writes 1500 bytes,
        forcing multiple \texttt{TtyTransmit} operations. The run exits with
        status 0 and prints \texttt{cp5\_tty: long write complete}.
\end{itemize}

\section*{Read Test}
\texttt{cp5\_read} verifies that a line can be split across two
\texttt{TtyRead} calls and that leftover bytes are returned by the next read.
The run was:

\begin{lstlisting}[basicstyle=\ttfamily\small]
printf "abcdef\nxyz\n" > cp5_input.txt
printf "\n" | xvfb-run -a timeout 8 ./yalnix -W -x -I0 cp5_input.txt cp5_read
\end{lstlisting}

The test exits with status 0. \texttt{TTYLOG.0} contains
\texttt{cp5\_read: ok}, and the trace shows two \texttt{TtyReceive} operations
for the two input lines.

\end{document}
